% Mục B.2. Phụ trách: TV3. Số RQ1 ở 3 triệu dòng do TV1 gửi (results/rq1_M1*.csv).
\phutrach{TV3 -- mục B.2}

\subsection{RQ1: thiết kế thí nghiệm}

\subsubsection{Bộ truy vấn}
\begin{table}[!htp]
\centering\bangnho
\caption{Bộ truy vấn vận đơn, thiết kế không covering và covering}\label{tab:truy-van}
\begin{tabular}{|C{1.1cm}|L{3.4cm}|L{4.1cm}|L{4.1cm}|C{1.6cm}|}
\hline
\dongdau \textbf{Mã} & \textbf{Nghiệp vụ} & \textbf{Không covering} & \textbf{B: covering} & \textbf{R} \\ \hline
Q1 & Tra một vận đơn & A: UNIQUE (tracking\_code); A2: cùng cột, không unique & (tracking\_code, status, updated\_at) & 1 \\ \hline
Q2 & 20 đơn mới nhất của một shop & A: (sender\_id, created\_at) & (sender\_id, created\_at, status, cod\_amount) & 20 \\ \hline
Q2R & Như Q2, bỏ LIMIT & Như Q2 & Như Q2 & 10 đến 10.000 \\ \hline
Q3 & Đếm đơn theo trạng thái tại một kho & A: (dest\_hub\_id) & (dest\_hub\_id, status) & khoảng N/200 \\ \hline
\end{tabular}
\end{table}

\subsubsection{Tách ba nguồn chênh lệch}
Một chênh lệch độ trễ giữa hai thiết kế chỉ mục có thể đến từ ba nguồn khác nhau, nên mỗi nguồn được đo riêng thay vì gán tất cả cho mạng.
\begin{itemize}
  \item \textbf{Độ chọn lọc} là số dòng $R$ phải đọc. Hai thiết kế trong một cặp so sánh luôn trả về cùng các dòng; Q2R chọn mốc thời gian để $R$ xấp xỉ 10, 100, 1.000 và 10.000.
  \item \textbf{Đường truy cập} là việc có phải quay về bảng hay không. Chỉ mục được ép bằng \code{FORCE INDEX} và kế hoạch được kiểm tra bằng \code{EXPLAIN} trước khi đo; nếu khác dự kiến thì dừng. Toán tử thực tế của mọi lượt nằm ở Bảng~\ref{tab:ke-hoach}.
  \item \textbf{Chi phí gọi KV} là phần cố định của mỗi pha truy cập, đo bằng số cop task, số RPC, thời gian xử lý và thời gian chờ trong \code{EXPLAIN ANALYZE} của TiDB. Phần còn lại là hiệu \code{rpc\_time} trừ xử lý trừ chờ; đó chỉ là xấp xỉ, gồm mã hóa, gRPC và lập lịch, không phải độ trễ mạng vì mọi node nằm trên một máy.
\end{itemize}
Hai đường chuẩn là \code{SELECT 1} (MySQL 0,362 ms, TiDB 0,571 ms) và tra theo khóa chính (MySQL 0,406 ms, TiDB 0,950 ms), đo trên M2 ở 3 triệu dòng. Bộ nhớ đệm kết quả của TiDB bị tắt trong lúc đo. Với Q1, thiết kế A là UNIQUE nên TiDB dùng \code{Point\_Get} thay vì coprocessor; vì vậy hiệu ứng covering của Q1 lấy từ A2 (cùng khóa nhưng không UNIQUE) so với B, còn A so với B trộn hai cơ chế khác nhau.

\subsection{RQ1: kế hoạch thực thi}
\begin{table}[!htp]
\centering\bangnho
\caption{Toán tử truy cập quan sát được}\label{tab:ke-hoach}
\begin{tabular}{|C{1.6cm}|C{2.4cm}|L{5.0cm}|L{5.0cm}|}
\hline
\dongdau \textbf{Truy vấn} & \textbf{Thiết kế} & \textbf{MySQL} & \textbf{TiDB C3} \\ \hline
Q1 & A / A2 / B & A: \code{IndexAndRow(const)}; A2: \code{IndexAndRow(ref)}; B: \code{CoveringIndex} & A: \code{Point\_Get}; A2: \code{IndexLookUp}; B: \code{IndexReader} \\ \hline
Q2, Q2R & A / B & A: \code{IndexAndRow(range)}; B: \code{CoveringIndex} & A: \code{IndexLookUp}; B: \code{IndexReader} \\ \hline
Q3 & A / B & A: \code{IndexAndRow(ref)}; B: \code{CoveringIndex} & A: \code{IndexLookUp}; B: \code{IndexReader} \\ \hline
\end{tabular}
\end{table}
Hình~\ref{fig:explain-tidb} và~\ref{fig:explain-mysql} chụp \code{EXPLAIN ANALYZE} của Q2 trên M2 (3 triệu dòng, một truy vấn mẫu, không phải số đo): trên TiDB, thiết kế A là \code{IndexLookUp} hai cop task (đọc chỉ mục rồi đọc dòng), thiết kế B là \code{IndexReader} một cop task; trên MySQL, thiết kế B là \code{Covering index range scan}.

\hinh{minh_chung_explain_tidb.png}{1.0}{Kế hoạch thực thi Q2 trên TiDB, thiết kế A (idx\_sender\_created) và B (idx\_sender\_cov); ảnh chụp terminal trên M2, một truy vấn mẫu}{fig:explain-tidb}

\hinh{minh_chung_explain_mysql.png}{1.0}{Kế hoạch thực thi Q2 trên MySQL, thiết kế A và B; ảnh chụp terminal trên M2, một truy vấn mẫu}{fig:explain-mysql}

Trong tệp \code{EXPLAIN ANALYZE} của lượt đo trên M3 (\path{results/raw}, lần lặp 1, 1 triệu dòng; truy vấn Q3 với \code{dest\_hub\_id=149}), một truy vấn mẫu cho thấy thêm:
\begin{itemize}
  \item \textbf{Q3, TiDB:} cả hai thiết kế đều quét chỉ mục 5\,003 khóa, nhưng thiết kế A còn phải đọc bảng qua 7 cop task (\code{tot\_proc} 253,3 ms, tổng \code{rpc\_info} 283 ms), còn thiết kế B gộp phép đếm xuống tầng chỉ mục với một cop task (\code{tot\_proc} 1,44 ms). Thời gian truy vấn mẫu là 88,3 ms so với 2,27 ms. Trên MySQL là 7,15 ms so với 0,608 ms: thiết kế A phải tạo bảng tạm để gom nhóm theo \code{status}, còn B đọc \code{status} thẳng từ chỉ mục và dùng \code{Group aggregate}.
\end{itemize}

\subsection{RQ1: lợi ích của covering}
Ở 3 triệu dòng trên M2 (số chính), lợi ích covering phụ thuộc vào $R$ và khác nhau giữa hai hệ.
\begin{itemize}
  \item \textbf{$R$ nhỏ} (Q1, Q2 và Q2R với $R = 10$): TiDB lợi hơn rõ, với tỉ số 1,48, 2,24 và 2,01; MySQL gần như không đổi, 0,99, 1,12 và 1,08.
  \item \textbf{$R = 100$}: 2,48 ở TiDB so với 1,56 ở MySQL.
  \item \textbf{$R \geq 1.000$}: MySQL lợi hơn, 3,56 so với 1,95 ở $R = 1.000$ và 5,91 so với 3,31 ở $R = 10.000$.
  \item \textbf{Q3} ($R = 15.000$, đếm theo trạng thái): hai hệ gần như bằng nhau, 16,78 ở MySQL và 16,86 ở TiDB.
  \item \textbf{Q1 với chỉ mục UNIQUE} (thiết kế A, \code{Point\_Get} trên TiDB): 1,608 ms, gần bằng thiết kế B (1,626 ms), nên tra một khóa duy nhất không cần covering.
\end{itemize}

\begin{table}[!htp]
\centering\bangnho
\caption[Độ trễ p50 và lợi ích covering trong từng hệ, 3 triệu dòng (M2)]{Độ trễ p50 và lợi ích covering $t(A)/t(B)$ trong từng hệ, 3 triệu dòng (M2)}\label{tab:covering}
\begin{tabular}{|C{1.5cm}|C{1.4cm}|C{1.8cm}|C{1.8cm}|C{1.8cm}|C{1.8cm}|C{1.6cm}|C{1.6cm}|}
\hline
\dongdau \textbf{Truy vấn} & \textbf{R} & \textbf{MySQL A (ms)} & \textbf{MySQL B (ms)} & \textbf{TiDB A (ms)} & \textbf{TiDB B (ms)} & \textbf{Lợi ích MySQL} & \textbf{Lợi ích TiDB} \\ \hline
Q1 (A2) & 1 & 0,429 & 0,434 & 2,408 & 1,626 & 0,99 & 1,48 \\ \hline
Q2 & 20 & 0,521 & 0,464 & 3,942 & 1,761 & 1,12 & 2,24 \\ \hline
Q2R & 10 & 0,498 & 0,462 & 3,718 & 1,847 & 1,08 & 2,01 \\ \hline
Q2R & 100 & 0,833 & 0,535 & 4,726 & 1,908 & 1,56 & 2,48 \\ \hline
Q2R & 1.000 & 3,309 & 0,929 & 7,791 & 3,987 & 3,56 & 1,95 \\ \hline
Q2R & 10.000 & 22,353 & 3,785 & 44,301 & 13,394 & 5,91 & 3,31 \\ \hline
Q3 & 15.000 & 28,266 & 1,685 & 111,876 & 6,634 & 16,78 & 16,86 \\ \hline
\end{tabular}
\par\vspace{3pt}{\noindent\small\itshape Nguồn: results/rq1\_M2\_summary.csv (số chính). Chỉ so lợi ích trong từng hệ.\par}
\end{table}

\hinh{rq1_M2_covering.png}{0.95}{Lợi ích covering theo truy vấn và $R$, 3 triệu dòng (M2)}{fig:covering}
Cùng phép đo ở 1 triệu dòng trên M3 (số liệu trong \path{results/rq1_M3_summary.csv}) cho cùng dạng với M2 ở các điểm có $R$ vừa và nhỏ: TiDB lợi hơn ở $R$ nhỏ (1,77 ở Q2 và 1,72 ở Q2R với $R = 10$, so với 1,10 và 1,07 của MySQL), còn MySQL lợi hơn từ $R = 1.000$ (2,41 so với 1,92). Ở $R = 10.000$ hai hệ gần nhau hơn trên M3 (4,03 ở MySQL, 3,80 ở TiDB) so với M2 (5,91 và 3,31). Q1 với chỉ mục không UNIQUE (A2) lợi 1,37 ở TiDB và 1,00 ở MySQL; Q1 với chỉ mục UNIQUE (\code{Point\_Get} trên TiDB) cho 0,94, tức không cần covering.

\textbf{Q3 trên TiDB không đáng tin để so.} Ở thiết kế A, p50 của 5 lượt là 99,4; 90,6; 102,8; 31,6 và 31,6 ms, và ngay trong một lượt, 5 truy vấn mẫu cho từ 15,4 đến 96,1 ms với cùng một kế hoạch; số lần đọc block ngoài bộ nhớ đệm bằng 0 ở cả năm lượt. Q2R với $R = 10.000$ ở thiết kế A cũng giảm theo cùng chiều (từ 91,1 xuống 42,8 ms). Nhóm không xác định được nguyên nhân. Tỉ số 24,6 phụ thuộc vào ba lượt chậm; nếu dùng hai lượt nhanh thì chỉ khoảng 8,6, gần bằng 8,6 của MySQL, nên báo cáo không kết luận hệ nào lợi hơn ở Q3 trên M3.

\subsection{RQ1: chi phí cố định, chi phí mỗi dòng và chi phí gọi KV}
\begin{table}[!htp]
\centering\bangfont
\caption{Mô hình $t = a + bR$ trên Q2R}\label{tab:mo-hinh}
\begin{tabular}{|C{2.4cm}|C{1.8cm}|C{1.8cm}|C{1.9cm}|C{2.2cm}|C{1.4cm}|C{2.0cm}|}
\hline
\dongdau \textbf{Máy, số dòng} & \textbf{Hệ} & \textbf{Thiết kế} & \textbf{a (ms)} & \textbf{b ($\mu$s/dòng)} & \textbf{$R^2$} & \textbf{Phần dư lớn nhất} \\ \hline
M2, 3 triệu & MySQL & A & 0,518 & 2,402 & 0,985 & 13\% \\ \hline
M2, 3 triệu & MySQL & B & 0,486 & 0,348 & 0,995 & 13\% \\ \hline
M2, 3 triệu & TiDB C3 & A & 3,918 & 3,977 & 0,999 & 14\% \\ \hline
M2, 3 triệu & TiDB C3 & B & 1,894 & 1,246 & 0,983 & 23\% \\ \hline
M1, 3 triệu & MySQL & A & 0,677 & 2,552 & 0,967 & 43\% \\ \hline
M1, 3 triệu & MySQL & B & 0,662 & 0,464 & 0,736 & 46\% \\ \hline
M1, 3 triệu & TiDB C3 & A & 7,476 & 7,668 & 0,907 & 45\% \\ \hline
M1, 3 triệu & TiDB C3 & B & 3,337 & 2,633 & 0,962 & 43\% \\ \hline
M3, 1 triệu & MySQL & A & 0,564 & 1,679 & 0,990 & 18\% \\ \hline
M3, 1 triệu & MySQL & B & 0,540 & 0,364 & 0,994 & 18\% \\ \hline
M3, 1 triệu & TiDB C3 & A & 3,095 & 4,619 & 0,818 & 46\% \\ \hline
M3, 1 triệu & TiDB C3 & B & 2,024 & 1,164 & 0,985 & 25\% \\ \hline
\end{tabular}
\par\vspace{3pt}{\noindent\small\itshape Khớp trên Q2R ($R$ từ 10 đến 10.000), trọng số $1/t$. Nguồn: results/rq1\_M2\_fit.csv, rq1\_M1\_fit.csv và rq1\_M3\_fit.csv. Không gộp các máy.\par}
\end{table}

Mô hình hai tham số giải thích dạng phụ thuộc vào $R$ này. Ở $R$ nhỏ, tỉ số tiến về $a_A / a_B$; ở $R$ lớn, tiến về $b_A / b_B$. Trên M2, $a_A / a_B$ là 1,07 ở MySQL và 2,07 ở TiDB; $b_A / b_B$ là 6,89 và 3,19. Nói cách khác, quay về bảng làm tăng chi phí \emph{cố định} của TiDB thêm 2,02 ms (MySQL chỉ thêm 0,033 ms), còn chi phí \emph{mỗi dòng} tăng thêm 2,73 $\mu$s ở TiDB và 2,05 $\mu$s ở MySQL. Chi phí mỗi dòng của đường covering trên TiDB (1,25 $\mu$s) cao hơn nhiều so với MySQL (0,35 $\mu$s), nên tỉ số tiệm cận của TiDB thấp hơn.

Các hệ số $a$, $b$ chỉ là phép khớp trong dải $R$ từ 10 đến 10.000 (độ lệch tương đối lớn nhất 23\% trên M2), không nên ngoại suy ra ngoài dải này.

\begin{table}[!htp]
\centering\bangnho
\caption{Chi phí gọi KV của TiDB theo truy vấn và thiết kế (3 triệu dòng, M2)}\label{tab:kv}
\begin{tabular}{|C{2.4cm}|C{1.3cm}|C{1.3cm}|C{1.4cm}|C{1.5cm}|C{1.7cm}|C{1.7cm}|C{2.0cm}|}
\hline
\dongdau \textbf{Truy vấn, R} & \textbf{Thiết kế} & \textbf{Cop task} & \textbf{Số RPC} & \textbf{Số key} & \textbf{Xử lý (ms)} & \textbf{Chờ (ms)} & \textbf{Phần còn lại (ms)} \\ \hline
Q1, R=1 & A2 & 2,0 & 2,0 & 2 & 0,16 & 0,11 & 0,72 \\ \hline
 & B & 1,0 & 1,0 & 1 & 0,07 & 0,06 & 0,33 \\ \hline
Q2, R=20 & A & 7,5 & 4,0 & 40 & 0,61 & 0,39 & 1,10 \\ \hline
 & B & 1,0 & 1,0 & 20 & 0,08 & 0,06 & 0,33 \\ \hline
Q2R, R=10 & A & 6,0 & 4,0 & 20 & 0,45 & 0,29 & 1,13 \\ \hline
 & B & 1,0 & 1,0 & 10 & 0,12 & 0,07 & 0,43 \\ \hline
Q2R, R=10.000 & A & 10,5 & 10,5 & 19.485 & 101,12 & 5,04 & 12,50 \\ \hline
 & B & 6,0 & 6,0 & 9.742 & 4,28 & 0,28 & 2,18 \\ \hline
Q3, R=15.000 & A & 18,0 & 18,0 & 30.079 & 465,57 & 257,73 & 22,93 \\ \hline
 & B & 1,0 & 1,0 & 15.040 & 4,06 & 0,07 & 0,48 \\ \hline
\end{tabular}
\par\vspace{3pt}{\noindent\small\itshape Trung vị của 10 lệnh \code{EXPLAIN ANALYZE} mỗi lượt. Nguồn: results/rq1\_M2\_summary.csv.\par}
\end{table}

Trên TiDB, thiết kế không covering cần nhiều cop task và RPC hơn: Q2 dùng 7,5 cop task và 4 RPC, trong khi thiết kế covering chỉ dùng 1 cop task và 1 RPC. Với $R = 10$, phần ``còn lại'' chiếm 58\% thời gian RPC của thiết kế A, còn với $R = 10.000$ chỉ 11\% vì thời gian xử lý ở TiKV chiếm phần lớn. Đây là bằng chứng cho chi phí cố định ở $R$ nhỏ và chi phí mỗi dòng ở $R$ lớn, nhưng phần còn lại chỉ là xấp xỉ nên không được viết là ``chi phí mạng''. Thời gian RPC và xử lý cộng dồn qua các tác vụ song song, nên có thể lớn hơn độ trễ đo được (ví dụ Q2R, $R = 10.000$, thiết kế A: RPC 118,8 ms so với p50 44,3 ms). Riêng Q3 thiết kế A, TiDB ghi nhận 1.086 lần đọc block ngoài bộ nhớ đệm trên M2 và 20 lần trên M1, trong khi \code{Innodb\_buffer\_pool\_reads} của MySQL bằng 0; vì vậy một phần chênh lệch ở Q3 có thể đến từ việc đọc đĩa chứ không chỉ từ chi phí gọi KV.

\subsection{RQ1: đối chiếu giữa các máy}
Phép đo lặp lại ở cùng quy mô trên M1 (\path{results/rq1_M1_summary.csv}) cho cùng một dạng: ở $R = 100$, tỉ số là 1,41 (MySQL) và 2,15 (TiDB), so với 1,56 và 2,48 trên M2; ở $R = 10.000$ là 5,65 và 2,83, so với 5,91 và 3,31. Khác biệt duy nhất là Q3: trên M1 MySQL lợi hơn (12,09 so với 9,11), còn trên M2 hai hệ bằng nhau, nên không kết luận hệ nào lợi hơn ở Q3. Trên M3 (1 triệu dòng), $a_A / a_B$ là 1,04 ở MySQL và 1,53 ở TiDB, $b_A / b_B$ là 4,62 và 3,97, cùng chiều với M2 (2,07 của TiDB và 6,89 của MySQL); phép khớp của TiDB thiết kế A kém ($R^2 = 0{,}82$) vì điểm $R = 10.000$ dao động 42,8--91,1 ms. Vì các máy khác nhau, chỉ có thể nói dạng kết quả lặp lại trên cả ba máy; độ trễ tuyệt đối không được so.

\subsection{RQ2: ngưỡng độ chọn lọc}

\subsubsection{Thiết kế}
Câu hỏi là: khi truy vấn lấy một khoảng giá trị của \code{cod\_amount}, tỉ lệ dòng khớp bao nhiêu thì việc tra qua chỉ mục \code{idx\_cod} rồi quay về bảng bắt đầu chậm hơn quét cả bảng, và optimizer của mỗi hệ có đổi phương án đúng chỗ đó hay không. Nhóm không đặt trước một điểm giao nhau chung cho hai hệ.

Truy vấn có dạng \code{SELECT id, tracking\_code, status, cod\_amount, created\_at FROM shipments WHERE cod\_amount BETWEEN lo AND hi}, chạy ở ba chế độ: \emph{index} (ép dùng \code{idx\_cod} bằng \code{FORCE INDEX}), \emph{scan} (loại chỉ mục bằng \code{IGNORE INDEX}) và \emph{auto} (không gợi ý, để optimizer chọn). Có 10 mức độ chọn lọc: 0,001; 0,01; 0,1; 0,5; 1; 2; 5; 10; 25 và 50\% số dòng. Khoảng \code{[lo, hi]} của từng mức được tính từ seed 20261116 theo phân phối thực của cột, nên số dòng khớp trung bình $R$ nằm sát mục tiêu, từ 30 dòng (0,001\%) đến khoảng 1,5 triệu dòng (50\%); cột $R$ ở Bảng~\ref{tab:rq2-p50} là số dòng thực đo. Từ mức 2\% trở xuống, mỗi lượt chạy 20 khoảng khác nhau; từ mức 5\% trở lên chạy 3 khoảng vì mỗi truy vấn đã lâu. Mỗi ô chạy 5 lượt, tổng $2 \times 10 \times 3 \times 5 = 300$ lượt trên M3, hồ sơ XS, bảng 3 triệu dòng; trước mỗi lượt chương trình kiểm tra kế hoạch thực thi (đúng toán tử chỉ mục hoặc quét bảng) và số dòng trả về giữa ba chế độ. Các giá trị trong bảng là trung vị p50 của 5 lượt. Chỉ so hai hệ theo hình dạng đường cong và điểm giao nhau; không so độ trễ tuyệt đối với M1 hoặc M2.

\subsubsection{Kết quả}
\begin{table}[!htp]
\centering\bangnho
\caption{Độ trễ p50 (ms) theo độ chọn lọc và chế độ, 3 triệu dòng (M3)}\label{tab:rq2-p50}
\begin{tabular}{|C{1.3cm}|C{1.7cm}|C{1.45cm}|C{1.45cm}|C{1.45cm}|C{1.45cm}|C{1.45cm}|C{1.45cm}|}
\hline
\dongdau & & \multicolumn{3}{c|}{\textbf{MySQL}} & \multicolumn{3}{c|}{\textbf{TiDB C3}} \\ \hline
\dongdau \textbf{Mức} & \textbf{$R$} & \textbf{index} & \textbf{scan} & \textbf{auto} & \textbf{index} & \textbf{scan} & \textbf{auto} \\ \hline
0,001\% & 30 & 0,656 & 386 & 0,785 & 4,32 & 882 & 3,89 \\ \hline
0,01\% & 304 & 1,44 & 383 & 1,48 & 4,49 & 887 & 4,47 \\ \hline
0,1\% & 3.011 & 7,66 & 393 & 7,65 & 18,9 & 901 & 17,5 \\ \hline
0,5\% & 15.053 & 36,6 & 371 & 36,2 & 82,8 & 909 & 71,0 \\ \hline
1\% & 29.987 & 76,6 & 435 & 76,0 & 198 & 911 & 184 \\ \hline
2\% & 59.935 & 154 & 453 & 150 & 426 & 939 & 931 \\ \hline
5\% & 150.029 & 368 & 489 & 374 & 2.097 & 981 & 980 \\ \hline
10\% & 300.183 & 752 & 540 & 756 & 4.888 & 1.055 & 1.061 \\ \hline
25\% & 749.587 & 1.898 & 722 & 529 & 10.874 & 866 & 842 \\ \hline
50\% & 1.500.344 & 3.810 & 981 & 1.021 & 26.352 & 1.792 & 1.756 \\ \hline
\end{tabular}
\par\vspace{3pt}{\noindent\small\itshape Nguồn: results/rq2\_M3\_summary.csv. Trung vị của 5 lượt; ô auto của TiDB ở mức 2\% có kế hoạch lai (một phần truy vấn dùng chỉ mục, một phần quét bảng).\par}
\end{table}

Ở độ chọn lọc thấp, chỉ mục thắng rất rõ: ở 0,001\% MySQL cần 0,66 ms (index) so với 386 ms (scan), TiDB cần 4,3 ms so với 882 ms. Chi phí quét bảng gần như không đổi theo $R$ (MySQL khoảng 0,37--0,54 s, TiDB khoảng 0,88--1,06 s cho tới 10\%), vì đây là chi phí đọc cả bảng; chi phí tra qua chỉ mục tăng theo $R$. Hai đường cắt nhau ở độ chọn lọc trung bình, và tăng $R$ thêm nữa thì quét bảng thắng.

\begin{table}[!htp]
\centering\bangfont
\caption{Điểm giao nhau và điểm optimizer đổi phương án}\label{tab:rq2}
\begin{tabular}{|C{2.2cm}|C{3.0cm}|C{4.2cm}|C{5.2cm}|}
\hline
\dongdau \textbf{Hệ} & \textbf{Điểm giao nhau (\%)} & \textbf{Optimizer chuyển sang quét bảng} & \textbf{Hậu quả ở mức đo gần nhất} \\ \hline
MySQL & 6,9 (giữa 5 và 10) & Giữa 10 và 25\%; ở 10\% vẫn dùng chỉ mục & Ở 10\%: auto 756 ms, scan 540 ms (auto chậm hơn 1,4 lần) \\ \hline
TiDB C3 & 3,2 (giữa 2 và 5) & Từ 2\% (kế hoạch lai), hoàn toàn ở 5\% & Ở 2\%: auto 931 ms, index 426 ms (auto chậm hơn 2,2 lần) \\ \hline
\end{tabular}
\par\vspace{3pt}{\noindent\small\itshape Nguồn: results/rq2\_M3\_summary.csv. Điểm giao nhau nội suy tuyến tính trong thang log giữa hai mức đo liên tiếp; vì lưới thưa nên chỉ khoảng trong ngoặc là chắc chắn.\par}
\end{table}

\textbf{Điểm giao nhau.} MySQL giao ở khoảng 6,9\%: ở 5\% index nhanh hơn scan (368 so với 489 ms), ở 10\% chậm hơn (752 so với 540 ms). TiDB giao sớm hơn, khoảng 3,2\%: ở 2\% index nhanh hơn (426 so với 939 ms), ở 5\% chậm hơn nhiều (2.097 so với 981 ms). Với cả hai hệ, khoảng giữa hai mức đo không chồng lên khoảng nhỏ nhất--lớn nhất của các lượt, nên việc điểm giao nằm trong khoảng đó không phụ thuộc vào nhiễu; con số 6,9\% và 3,2\% chỉ là nội suy trong khoảng ấy.

\textbf{Optimizer.} Hai hệ lệch khỏi điểm giao nhau theo hai chiều ngược nhau. MySQL giữ chỉ mục quá lâu: ở 10\% optimizer vẫn chọn chỉ mục dù quét bảng nhanh hơn, và chỉ đổi sang quét bảng ở mức 25\% (đo được trong khoảng từ 10 đến 25\%). TiDB đổi quá sớm: ở 2\%, mức mà chỉ mục còn nhanh gấp 2,2 lần, kế hoạch \code{auto} đã lẫn \code{IndexLookUp} với \code{TableReader} ở cả 5 lượt, và p50 của \code{auto} (931 ms) gần bằng \code{scan} (939 ms) thay vì \code{index} (426 ms). Khoảng nhỏ nhất--lớn nhất của \code{auto} và \code{index} ở mức này không chồng nhau (706--941 so với 412--595 ms). Số liệu không cho biết nguyên nhân: nhóm không đo thống kê mà optimizer dùng để ước lượng số dòng, nên không kết luận optimizer ước lượng sai chỗ nào.

\textbf{Chi phí cố định và chi phí trên mỗi dòng.} Khớp mô hình $t = a + bR$ với 10 mức (trọng số $1/t$, \path{results/rq2_M3_fit.csv}) cho: MySQL index $a = 0{,}60$ ms, $b = 2{,}49$ $\mu$s/dòng; MySQL scan $a = 397$ ms, $b = 0{,}42$ $\mu$s/dòng, R$^2$ = 0,98; TiDB index $a = 3{,}05$ ms, $b = 7{,}17$ $\mu$s/dòng; TiDB scan $a = 890$ ms, $b = 0{,}30$ $\mu$s/dòng. Với MySQL, mô hình khớp tốt (R$^2$ = 0,9995 cho index): chi phí khoảng 2,5 $\mu$s mỗi dòng, hầu như không đổi trên toàn dải $R$ (từ 2,4 đến 2,6 $\mu$s). Với TiDB, mô hình tuyến tính khớp kém (R$^2$ = 0,56 cho index và 0,60 cho scan, phần dư lớn nhất 59\%). Lý do đọc được từ số liệu là chi phí trên mỗi dòng của TiDB không cố định: thời gian chia cho $R$ tăng từ khoảng 6,6 $\mu$s ở 1\% lên 14--18 $\mu$s từ 5\% trở lên. Vì vậy hệ số $b$ của TiDB (7,2 $\mu$s) chỉ là giá trị trung bình thô, và bảng này không được dùng để dự đoán $R$ nằm ngoài các mức đã đo. Nhóm chưa tách được nguyên nhân của sự tăng này; các bộ đếm chi tiết được lưu trong \path{results/rq2_M3.csv}.

\hinh{rq2_M3.png}{0.95}{Độ trễ p50 theo độ chọn lọc, thang log cả hai trục, ba chế độ; đường dọc là điểm giao nhau của từng hệ (M3, 3 triệu dòng)}{fig:rq2}

\textbf{Độ tin cậy.} Quét bảng của TiDB dao động nhiều giữa các lượt (656--1.139 ms ở cùng một mức, so với 359--389 ms của MySQL). Ở mức 25\%, hai chế độ của MySQL chạy cùng kế hoạch quét bảng nhưng cho 722 và 529 ms, nên chênh lệch dưới khoảng một phần ba được coi là nhiễu và kết luận chỉ dựa trên các khác biệt lớn hơn. Điểm giao nhau có thể dịch chuyển khi dữ liệu nhỏ hơn hoặc lớn hơn bộ nhớ đệm.

\subsection{Kết luận chương}
\textbf{RQ1: covering index và chi phí đọc.} Giả thuyết được xác nhận một phần. Ở $R$ nhỏ, covering lợi cho TiDB hơn (Q2R, $R = 10$: 2,01 so với 1,08 trên M2; 1,72 so với 1,07 trên M3), vì quay về bảng thêm vào TiDB khoảng 2,0 ms cố định so với 0,03 ms ở MySQL. Ở $R$ lớn, hai hệ không gần nhau mà MySQL lợi hơn trên M2 (5,91 so với 3,31 khi $R = 10.000$; trên M3 là 4,03 so với 3,80), có thể vì mỗi dòng quay về bảng ở MySQL là một lần tra B+Tree; điều này chưa được kiểm chứng. Ở Q3 không kết luận được hệ nào lợi hơn.

\textbf{RQ2: ngưỡng độ chọn lọc.} Điểm giao nhau giữa dùng chỉ mục và quét bảng là khoảng 6,9\% ở MySQL và 3,2\% ở TiDB (M3, 3 triệu dòng). Optimizer lệch theo hai chiều ngược nhau: MySQL giữ chỉ mục quá lâu (ở 10\%, chậm hơn quét bảng 1,4 lần), còn TiDB bỏ chỉ mục quá sớm (ở 2\%, chậm hơn dùng chỉ mục 2,2 lần).

